% ==================== 页面与字体 ====================
% 使用 TeX Live 自带 Fandol 字体，避免依赖本机 Windows 字体。
\usepackage[a4paper,top=2.5cm,bottom=2.5cm,left=2.5cm,right=2.5cm]{geometry}
\usepackage{setspace}
\usepackage{indentfirst}
\setstretch{1.28}
\setlength{\parindent}{2em}
\setlength{\parskip}{0pt}

% ==================== 数学、表格与图片 ====================
\usepackage{amsmath,amssymb,mathtools,bm}
\usepackage{booktabs,multirow,array,longtable,tabularx,pdflscape}
\usepackage{graphicx,float,subcaption,standalone,svg}
\usepackage{tikz}
\usetikzlibrary{arrows.meta,positioning}
\usepackage{pgfplots}
\pgfplotsset{compat=1.18}

% 主文档默认由 paper/ 作为工作目录编译，因此数据与结果目录位于上一级。
\newcommand{\RawDataDir}{../data/raw}
\newcommand{\ProcessedDataDir}{../data/processed}
\newcommand{\OutputsDir}{../outputs}
\newcommand{\QOneResultsDir}{../outputs/q1/tables}
\newcommand{\QTwoResultsDir}{../outputs/q2/tables}
\newcommand{\QThreeResultsDir}{../outputs/q3/tables}

\usepackage{siunitx}
\sisetup{detect-all}

% ==================== 版式与引用 ====================
\usepackage{enumitem}
\usepackage{titlesec}
\usepackage{fancyhdr}
\usepackage{xcolor}
\usepackage{hyperref}
\hypersetup{hidelinks,pdfauthor={},pdftitle={\PaperTitle}}

% 页码位于页脚中部；正文与附录均不放身份信息。
\fancypagestyle{cumcm}{%
  \fancyhf{}
  \fancyfoot[C]{\thepage}
  \renewcommand{\headrulewidth}{0pt}
  \renewcommand{\footrulewidth}{0pt}
}

% 标题层级：不使用目录，但保留清晰层次。
% 使用非星号 \titlespacing，使标题后的首段保留正常的 2em 首行缩进。
\titleformat{\section}{\centering\heiti\zihao{3}}{\thesection}{1em}{}
\titleformat{\subsection}{\heiti\zihao{4}}{\thesubsection}{0.8em}{}
\titleformat{\subsubsection}{\heiti\zihao{-4}}{\thesubsubsection}{0.6em}{}
\titlespacing{\section}{0pt}{1.4em}{0.8em}
\titlespacing{\subsection}{0pt}{1.0em}{0.5em}
\titlespacing{\subsubsection}{0pt}{0.8em}{0.3em}

\setlist[itemize]{leftmargin=2.5em,itemsep=0.2em,topsep=0.3em}
\setlist[enumerate]{leftmargin=2.8em,itemsep=0.2em,topsep=0.3em}

% 常用列类型。
\newcolumntype{Y}{>{\centering\arraybackslash}X}
\newcolumntype{L}[1]{>{\raggedright\arraybackslash}p{#1}}
\newcolumntype{C}[1]{>{\centering\arraybackslash}p{#1}}

% 骨架中的待填文字；交稿前应全部替换或删除。
% 使用 parbox 允许较长提示换行，参数中的特殊字符仍需按 LaTeX 规则转义。
\newcommand{\todo}[1]{\textcolor{gray}{\parbox{\linewidth}{【待补：#1】}}}
\newcommand{\modelname}[1]{\textbf{#1}}

\renewcommand{\figurename}{图}
\renewcommand{\tablename}{表}
\numberwithin{equation}{section}
\numberwithin{figure}{section}
\numberwithin{table}{section}
